Assume that $A$ always outputs a hypothesis consistent with $D$. Occam's Razor suggests:
\begin{center}
    among all consistent hypotheses, choose the \textbf{simplest}.
\end{center}
\textbf{Simplicity} is formalized mathematically in two distinct ways:
\begin{enumerate}
    \item \textbf{Description length}, how many bits it takes to write it down using our representation language.
    \item \textbf{Size} of the hypothesis, number of hyper-parameters required.
\end{enumerate} 
The take-away message must be:
\begin{center}
    fewer simple hypotheses $\rightarrow$ less chance of fitting noise, so \textbf{overfitting}.
\end{center}
Simple hypotheses guarantee \textbf{better generalization behaviors}.